% ============================================================================
% 第三章：基于表格型 Q-Learning 的时间窗优化方法
% 优化说明：修复页面超出问题，添加奖励计算和预测流程
\section{方法}
% ============================================================================

{\small

\subsection{问题形式化}

\begin{frame}{问题形式化：MDP 建模}
    \begin{block}{马尔可夫决策过程五元组 $(\mathcal{S}, \mathcal{A}, \mathcal{R}, \mathcal{P}, \gamma)$}
        \begin{itemize}
            \item \textbf{状态空间 $\mathcal{S}$}\cite{blankertz2008optimizing}：$(t_{\text{start}}, t_{\text{end}})$
            \begin{itemize}
                \item 步长$\Delta t = 0.2$秒，$t_{\text{end}} - t_{\text{start}} \geq 0.5$秒
                \item 有效状态：\textbf{136 个}
            \end{itemize}
            \item \textbf{动作空间 $\mathcal{A}$}：$\{a_0, a_1, a_2, a_3\}$（边界调整）
            \item \textbf{奖励函数 $\mathcal{R}$}：$r(s, a) = \text{Accuracy}_{\text{CV}}(s')$
            \item \textbf{状态转移 $\mathcal{P}$}：确定性过程
        \end{itemize}
    \end{block}

    \vspace{0.2cm}
    \begin{alertblock}{优化目标}
        \centering
        $\pi^* = \arg\max_{\pi} \mathbb{E}_{\pi} \left[ \sum_{t=0}^{\infty} \gamma^t r_t \right]$
    \end{alertblock}
\end{frame}

% ============================================================================
% 表格型 Q-Learning 算法
% ============================================================================
\subsection{表格型 Q-Learning 算法}

\begin{frame}{Q-Learning 算法原理}
    \begin{columns}[T]
        \column{0.48\textwidth}
        \begin{block}{Q 函数定义\cite{watkins1992q}}
            \footnotesize
            $Q(s, a)$ 表示状态$s$执行动作$a$后的期望累积奖励：
            \begin{equation*}
            \footnotesize
            Q(s, a) = \mathbb{E}\left[\sum_{t=0}^{\infty} \gamma^t r_t\right]
            \end{equation*}
            最优 Q 函数满足 Bellman 方程：
            \begin{equation*}
            \footnotesize
            Q^*(s, a) = \mathbb{E}\left[ r + \gamma \max_{a'} Q^*(s', a') \right]
            \end{equation*}
        \end{block}

        \column{0.48\textwidth}
        \begin{block}{时序差分更新}
            \footnotesize
            Q 值更新规则：
            \begin{equation*}
            \footnotesize
            Q(s_t, a_t) \leftarrow Q(s_t, a_t) + \alpha \cdot \text{TD}_t
            \end{equation*}
            其中 TD 误差：
            \begin{equation*}
            \footnotesize
            \text{TD}_t = r_{t+1} + \gamma \max_{a} Q(s_{t+1}, a) - Q(s_t, a_t)
            \end{equation*}
        \end{block}
    \end{columns}

    \begin{exampleblock}{核心思想}
        \footnotesize
        \begin{itemize}
            \item 维护 Q 表，无需函数近似
            \item 通过 TD 误差迭代更新，自举学习
            \item $\varepsilon$-贪婪策略平衡探索与利用
        \end{itemize}
    \end{exampleblock}
\end{frame}

\begin{frame}{Q-Learning 时间窗优化算法伪代码}
    \begin{algorithm}[H]
    \scriptsize
    \begin{algorithmic}[1]
    \REQUIRE EEG 训练数据$\mathcal{D}$，状态空间$\mathcal{S}$，动作空间$\mathcal{A}$
    \ENSURE 最优时间窗$(t_{\text{start}}^*, t_{\text{end}}^*)$

    \STATE 初始化$Q[s, a] \leftarrow 0, \forall s, a$
    \FOR{$\text{episode} = 1$ \TO $N$}
        \STATE 随机初始化状态$s_0$
        \FOR{$t = 0$ \TO $T-1$}
            \IF{$\text{random} < \varepsilon$}
                \STATE 随机选择动作$a_t$
            \ELSE
                \STATE $a_t \leftarrow \arg\max_a Q[s_t, a]$
            \ENDIF
            \STATE 执行动作：$s_{t+1} \leftarrow f(s_t, a_t)$
            \STATE 计算奖励：$r_t \leftarrow \text{EvaluateWindow}(s_{t+1}, \mathcal{D})$
            \STATE 更新 Q 值：$Q[s_t, a_t] \leftarrow Q[s_t, a_t] + \alpha[r_t + \gamma\max_a Q[s_{t+1}, a] - Q[s_t, a_t]]$
            \STATE $s_t \leftarrow s_{t+1}$
        \ENDFOR
        \STATE $\varepsilon \leftarrow \max(\varepsilon_{\text{min}}, \varepsilon \times \text{decay})$
    \ENDFOR
    \STATE \RETURN $s^* = \arg\max_s \max_a Q[s, a]$
    \end{algorithmic}
    \end{algorithm}
\end{frame}

\begin{frame}{算法流程图}
    \begin{center}
    \resizebox{1.0\textwidth}{!}{%
    \begin{tikzpicture}[
        node distance=0.4cm,
        block/.style={rectangle, draw, rounded corners, minimum width=2cm, minimum height=0.45cm, align=center},
        decision/.style={diamond, draw, aspect=2, minimum width=1.8cm, minimum height=0.45cm, align=center},
        arrow/.style={->, thick}
    ]
    % 第一行
    \node[block, fill=blue!10] (init) {初始化 Q 表};
    \node[block, fill=green!10, right=of init] (episode) {开始 Episode};
    \node[block, fill=yellow!10, right=of episode] (action) {$\varepsilon$-贪婪\\选择动作};
    
    % 第二行
    \node[block, fill=orange!10, below=of action] (transition) {状态转移};
    \node[block, fill=red!10, left=of transition] (reward) {计算奖励};
    \node[block, fill=purple!10, right=of transition] (update) {更新 Q 值};
    
    % 第三行
    \node[decision, below=of transition] (check) {达到最大\\步数？};
    \node[block, fill=gray!10, below=of check] (end) {返回最优$s^*$};
    
    % 连接箭头
    \draw[arrow] (init) -- (episode);
    \draw[arrow] (episode) -- (action);
    \draw[arrow] (action) -- (transition);
    \draw[arrow] (transition) -- (reward);
    \draw[arrow] (transition) -- (update);
    \draw[arrow] (update) |- (check);
    \draw[arrow] (reward) |- (check);
    \draw[arrow] (check.east) -- ++(0.6,0) |- (action);
    \draw[arrow] (check) -- (end);
    \end{tikzpicture}%
    }
    \end{center}
\end{frame}

% ============================================================================
% 奖励计算与预测
% ============================================================================
\subsection{奖励计算与预测}

\begin{frame}{奖励计算：EvaluateWindow}
    \begin{block}{输入与输出}
        \footnotesize
        \begin{itemize}
            \item \textbf{输入}：状态$s = (t_{\text{start}}, t_{\text{end}})$，EEG 训练数据$\mathcal{D}$
            \item \textbf{输出}：5 折交叉验证准确率$\text{acc}$
        \end{itemize}
    \end{block}

    \vspace{0.2cm}
    \begin{alertblock}{计算步骤\cite{blankertz2008optimizing}}
        \footnotesize
        \begin{enumerate}
            \item \textbf{截取时间窗}：$\mathbf{X}_i^{\text{window}} \leftarrow \mathbf{X}_i[:, t_{\text{start}}:t_{\text{end}}]$
            \item \textbf{多类别 CSP 特征提取}：
            \begin{itemize}
                \item 计算各类别协方差矩阵：$\mathbf{C}_c = \frac{1}{N_c}\sum_{i \in \text{class}_c} \mathbf{X}_i^{\text{window}} (\mathbf{X}_i^{\text{window}})^\top$
                \item 近似对角化：$\mathbf{C}_1 \mathbf{w} = \lambda (\mathbf{C}_1 + \mathbf{C}_2 + \mathbf{C}_3 + \mathbf{C}_4) \mathbf{w}$
                \item 构建空间滤波器$\mathbf{W}$（4 个分量）
            \end{itemize}
            \item \textbf{对数方差特征}：$\mathbf{f}_i = \log(\text{var}(\mathbf{W} \mathbf{X}_i^{\text{window}}))$
            \item \textbf{Z-score 标准化}：$\mathbf{f}_i \leftarrow (\mathbf{f}_i - \mu_f) / \sigma_f$
            \item \textbf{5 折交叉验证}：训练线性 SVM（4 类分类），计算平均准确率
        \end{enumerate}
    \end{alertblock}
\end{frame}

\begin{frame}{预测流程}
    \begin{columns}[T]
        \column{0.48\textwidth}
        \begin{block}{训练阶段（离线）}
            \footnotesize
            \begin{enumerate}
                \item 初始化 Q 表
                \item 执行 Q-Learning 算法
                \item 每步调用 EvaluateWindow
                \item 迭代更新 Q 值
                \item 输出最优时间窗$s^*$
            \end{enumerate}
            
            \vspace{0.2cm}
            \textbf{计算开销}：4000 次奖励评估（一次性）
        \end{block}

        \column{0.48\textwidth}
        \begin{alertblock}{推理阶段（在线）}
            \footnotesize
            \begin{enumerate}
                \item 加载训练好的 Q 表
                \item 查表找到最优状态：
                \begin{equation*}
                \footnotesize
                s^* = \arg\max_{s} \max_{a} Q(s, a)
                \end{equation*}
                \item 使用$s^*$截取 EEG 信号
                \item CSP-SVM 分类预测
            \end{enumerate}
            
            \vspace{0.2cm}
            \textbf{时间复杂度}：$O(544) \approx O(1)$
        \end{alertblock}
    \end{columns}

    \vspace{0.2cm}
    \begin{exampleblock}{效率对比}
        \footnotesize
        \begin{tabular}{l|c|c}
            \hline
            \textbf{阶段} & \textbf{表格型 Q-Learning} & \textbf{DQN}\cite{mnih2015human} \\
            \hline
            训练 & 4000 次评估 & 百万级参数训练 \\
            推理 & $O(1)$查表 & $O(\text{网络前向传播})$ \\
            \hline
        \end{tabular}
    \end{exampleblock}
\end{frame}

% ============================================================================
% 收敛性分析
% ============================================================================
\subsection{收敛性分析}

\begin{frame}{收敛性保证}
    \begin{block}{Q-Learning 收敛性定理\cite{watkins1992q}}
        \footnotesize
        对于有限 MDP，若满足：
        \begin{itemize}
            \item 所有状态 - 动作对被无限次访问
            \item 学习率序列满足 Robbins-Monro 条件
            \item 采用适当探索策略（如$\varepsilon$-贪婪）
        \end{itemize}
        则 Q 值序列以概率 1 收敛到最优动作价值函数$Q^*$
    \end{block}

    \begin{alertblock}{本研究的收敛性条件}
        \footnotesize
        \begin{itemize}
            \item \textbf{有限状态空间}：136 状态 × 4 动作 = 544 个 Q 值条目
            \item \textbf{充分探索}：$\varepsilon$-贪婪策略，4000 步充分访问
            \item \textbf{确定性转移}：时间窗参数调整是可预测的离散操作
        \end{itemize}
    \end{alertblock}

    \begin{exampleblock}{结论}
        \footnotesize
        本研究算法在理论上保证收敛到最优时间窗选择策略
    \end{exampleblock}
\end{frame}

% ============================================================================
% 与 DQN 的理论对比
% ============================================================================
\subsection{与 DQN 的理论对比}

\begin{frame}{表格型 Q-Learning vs DQN}
    \begin{center}
    \footnotesize
    \renewcommand{\arraystretch}{1.4}
    \begin{tabular}{l|cc}
        \hline
        \textbf{维度} & \textbf{表格型 Q-Learning} & \textbf{DQN}\cite{mnih2015human} \\
        \hline
        价值函数表示 & $Q(s,a)$表格存储 & $Q(s,a;\theta)$神经网络近似 \\
        更新规则 & $Q \leftarrow Q + \alpha \cdot \text{TD}$ & $\theta \leftarrow \theta - \eta \nabla_\theta \mathcal{L}$ \\
        收敛性保证 & 理论保证\cite{watkins1992q} & 经验性收敛 \\
        样本效率 & 单步传播奖励 & 批量梯度下降 \\
        推理复杂度 & $O(1)$查表 & $O(\text{网络前向传播})$ \\
        可解释性 & Q 值透明可视 & 黑盒决策 \\
        \hline
    \end{tabular}
    \end{center}

    \vspace{0.2cm}
    \begin{alertblock}{DQN 损失函数}
        \footnotesize
        \begin{equation*}
        \mathcal{L}(\theta) = \mathbb{E}\left[ \left( r + \gamma \max_{a'} Q(s', a'; \theta^-) - Q(s, a; \theta) \right)^2 \right]
        \end{equation*}
    \end{alertblock}
\end{frame}

\begin{frame}{方法选择依据}
    \begin{columns}[T]
        \column{0.48\textwidth}
        \begin{block}{本任务特性}
            \footnotesize
            \begin{itemize}
                \item \textbf{状态空间规模}：$|\mathcal{S}| = 136$（小规模）
                \item \textbf{状态可离散化}：时间窗参数精确离散
                \item \textbf{转移确定性}：参数调整可预测
                \item \textbf{数据稀缺}：BCI 试次有限
            \end{itemize}
        \end{block}

        \column{0.48\textwidth}
        \begin{alertblock}{选择表格型的理由}
            \footnotesize
            \begin{itemize}
                \item Q 表仅 544 条目（约 4 KB 内存）
                \item 严格收敛性保证
                \item 单步奖励传播，样本效率高
                \item Q 值表可直接可视化分析
            \end{itemize}
        \end{alertblock}
    \end{columns}

    \vspace{0.2cm}
    \begin{exampleblock}{核心结论}
        \footnotesize
        在时间窗优化（136 状态）任务中，表格型 Q-Learning 是更合适的选择
    \end{exampleblock}
\end{frame}

% ============================================================================
% 本章小结
% ============================================================================
\subsection{本章小结}

\begin{frame}{本章小结}
    \begin{block}{方法创新点}
        \footnotesize
        \begin{itemize}
            \item \textbf{问题建模简洁性}：离散化避免深度强化学习复杂性
            \item \textbf{理论保证严格性}：有限状态空间下收敛性保证
            \item \textbf{计算效率优势}：离线训练 4000 次评估，在线推理$O(1)$查表
            \item \textbf{可解释性强}：Q 值表可直接可视化和分析
        \end{itemize}
    \end{block}

    \vspace{0.2cm}
    \begin{alertblock}{局限性}
        \footnotesize
        \begin{itemize}
            \item 状态空间需人工离散化，可能丢失连续信息
            \item 奖励计算依赖 5 折交叉验证，训练开销较大
            \item 仅适用于时间窗优化特定任务
        \end{itemize}
    \end{alertblock}
\end{frame}

} % end of \small
